% Banque d'exercices — chapitre 4
% Le chapitre est déterminé par le nom de ce fichier.

\begin{exo}[Travail et chaleur le long de trois chemins élémentaires]{travail-trois-chemins}
\keywords{travail des forces de pression, isotherme, isobare, isochore, diagramme de Clapeyron, premier principe}

On considère un système fermé constitué de $n$ moles de gaz parfait monoatomique ($U = \frac{3}{2}nRT$), initialement dans l'état $A(P_A, V_A, T_A)$. Toutes les transformations sont quasi-statiques.

\begin{enumerate}
\item Compression isotherme $A\to B$ jusqu'à $V_B = V_A/2$. Calculer $P_B$, puis $W$, $\Delta U$ et $Q$ sur cette branche.
\item Détente isobare $B\to C$ jusqu'à $V_C = V_A$. Calculer $T_C$, puis $W$, $\Delta U$ et $Q$.
\item Refroidissement isochore $C\to A$ qui referme le cycle. Calculer $W$, $\Delta U$ et $Q$.
\item Représenter le cycle $A\to B\to C\to A$ dans le diagramme de Clapeyron $(V,P)$. Application numérique pour $n = 1$ mol, $P_A = 1{,}0$ bar, $T_A = 300$ K : donner $(P,V,T)$ en chacun des trois points.
\item Comparer $W_{A\to B}$ au travail $W_{A\to C\to B}$ reçu le long du chemin isochore $A\to C$ puis isobare $C\to B$. Commenter.
\item Le cycle $A\to B\to C\to A$ est-il moteur ou récepteur ? Calculer le travail total et vérifier la cohérence avec le sens de parcours.
\item Représenter le même cycle dans le diagramme $(T,P)$.
\end{enumerate}

\begin{indication}
Sur une isotherme de gaz parfait, remplacer $P$ par $nRT/V$ dans $\delta W = -P\,\mathrm{d}V$ avant d'intégrer. En diagramme $(V,P)$, l'aire du cycle vaut le travail fourni si le parcours est horaire. Convention banquier : tout ce qui entre dans le système est compté positivement.
\end{indication}

\begin{solution}
\textbf{Question 1.} Isotherme : $T_B = T_A$, et $P_AV_A = P_BV_B$ donne $P_B = 2P_A$.
\[
W_{AB} = -\int_{V_A}^{V_A/2}\frac{nRT_A}{V}\,\mathrm{d}V = -nRT_A\ln\frac{V_B}{V_A} = nRT_A\ln 2 > 0.
\]
$U$ ne dépend que de $T$ : $\Delta U = 0$, donc $Q_{AB} = -W_{AB} = -nRT_A\ln 2 < 0$ (le gaz évacue la chaleur de compression).

\textbf{Question 2.} $P_CV_C = 2P_AV_A = 2nRT_A$ donc $T_C = 2T_A$. À pression constante $P_B = 2P_A$ :
\[
W_{BC} = -2P_A\left(V_A - \tfrac{V_A}{2}\right) = -P_AV_A = -nRT_A, \qquad \Delta U_{BC} = \tfrac{3}{2}nR(T_C - T_A) = \tfrac{3}{2}nRT_A,
\]
d'où $Q_{BC} = \Delta U_{BC} - W_{BC} = \tfrac{5}{2}nRT_A$.

\textbf{Question 3.} Isochore : $W_{CA} = 0$ et $Q_{CA} = \Delta U_{CA} = \tfrac{3}{2}nR(T_A - 2T_A) = -\tfrac{3}{2}nRT_A$.

\textbf{Question 4.} $V_A = RT_A/P_A \approx 24{,}9$ L. Donc $A$ : $(1$ bar, $24{,}9$ L, $300$ K$)$, $B$ : $(2$ bar, $12{,}5$ L, $300$ K$)$, $C$ : $(2$ bar, $24{,}9$ L, $600$ K$)$.

\textbf{Question 5.} $W_{A\to C\to B} = 0 + \left[-2P_A\left(\tfrac{V_A}{2} - V_A\right)\right] = P_AV_A = nRT_A \approx 2494$ J, alors que $W_{A\to B} = nRT_A\ln 2 \approx 1729$ J. Mêmes états extrêmes, travaux différents : le travail dépend du chemin suivi, $\delta W$ n'est pas une différentielle exacte.

\textbf{Question 6.} $W_{\text{total}} = nRT_A\ln 2 - nRT_A = nRT_A(\ln 2 - 1) \approx -765$ J $< 0$ : le gaz fournit du travail, le cycle est moteur. C'est cohérent : en diagramme $(V,P)$, le parcours est horaire (isotherme en bas, isobare en haut) et l'aire du cycle vaut le travail fourni.

\textbf{Question 7.} En diagramme $(T,P)$ : $A\to B$ est un segment vertical à $T = T_A$ (de $P_A$ à $2P_A$) ; $B\to C$ un segment horizontal à $P = 2P_A$ (de $T_A$ à $2T_A$) ; $C\to A$ un segment de la droite passant par l'origine $P = (nR/V_A)\,T$ (isochore).
\end{solution}
\end{exo}

\begin{exo}[Chauffage à pression constante : découverte de $C_P$]{decouverte-cp}
\keywords{capacité thermique, pression constante, relation de Mayer, chauffage isobare, gaz parfait}

Une mole de gaz parfait monoatomique ($C_V = \frac{3}{2}R$) est enfermée dans une enceinte de volume initial $V_0 = 25$ L, fermée par un piston mobile sans frottement soumis à la pression atmosphérique $P_0 = 1{,}0\times10^5$ Pa, supposée constante. Le système est initialement à l'équilibre. Une petite résistance électrique placée dans l'enceinte apporte une chaleur $Q = 1000$ J de façon quasi-statique.

\begin{enumerate}
\item Déterminer numériquement la température initiale $T_0$.
\item Déterminer l'état final (température, pression, volume) de façon littérale puis numérique. Calculer $W$ et $\Delta U$. Quelle fraction de la chaleur apportée sert à repousser l'atmosphère ?
\item Ce calcul permet de définir la capacité thermique molaire à pression constante, $C_P = \left(\delta Q/\mathrm{d}T\right)_P$. Exprimer $C_P$ en fonction de $C_V$ et $R$. Cette formule est un cas particulier d'une relation générale que l'on démontrera plus tard (relation de Mayer).
\end{enumerate}

\begin{indication}
La transformation est isobare : $W = -P_0\,\Delta V$, et l'équation d'état donne $P_0\,\Delta V = R\,\Delta T$. Injecter dans le premier principe pour obtenir une équation sur $\Delta T$ seul.
\end{indication}

\begin{solution}
\textbf{Question 1.} $T_0 = \dfrac{P_0V_0}{R} = \dfrac{10^5 \times 25\times10^{-3}}{8{,}314} \approx 301$ K.

\textbf{Question 2.} La pression finale vaut $P_0$ (équilibre mécanique du piston). Premier principe :
\[
\Delta U = Q + W = Q - P_0\,\Delta V.
\]
Or $\Delta U = \tfrac{3}{2}R\,\Delta T$ et $P_0\,\Delta V = R\,\Delta T$ (isobare), donc
\[
\tfrac{3}{2}R\,\Delta T = Q - R\,\Delta T \implies \Delta T = \frac{2Q}{5R} \approx 48\ \text{K}.
\]
D'où $T_f \approx 349$ K et $V_f = V_0\,T_f/T_0 \approx 29{,}0$ L. On en tire les valeurs exactes
\[
W = -P_0\,\Delta V = -\tfrac{2}{5}Q = -400\ \text{J}, \qquad \Delta U = \tfrac{3}{5}Q = +600\ \text{J}.
\]
Les $2/5$ de la chaleur fournie repartent en travail pour repousser l'atmosphère ; seuls les $3/5$ restants chauffent le gaz.

\textbf{Question 3.} À pression constante, $\delta Q = \mathrm{d}U + P\,\mathrm{d}V = C_V\,\mathrm{d}T + R\,\mathrm{d}T$, donc
\[
C_P = C_V + R = \tfrac{5}{2}R \approx 20{,}8\ \text{J·K}^{-1}\text{mol}^{-1}.
\]
On retrouve bien $Q/\Delta T = 1000/48{,}1 \approx 20{,}8$ J/K. $C_P > C_V$ : il faut plus de chaleur pour élever la température à pression constante, car une partie de l'énergie s'échappe en travail.
\end{solution}
\end{exo}

\begin{exo}[Détente de Joule-Gay-Lussac : irréversibilité et dépendance au chemin]{detente-joule-gay-lussac}
\keywords{detente dans le vide, transformation irreversible, differentielle inexacte, gaz parfait, premier principe}

Un récipient rigide et calorifugé est séparé en deux compartiments de même volume $V_0$ par une paroi amovible. Le compartiment de gauche contient $n$ moles de gaz parfait dans l'état $(P_0, V_0, T_0)$ ; celui de droite est vide. On retire brusquement la paroi (détente dite de Joule--Gay-Lussac, étudiée par Gay-Lussac en 1807 puis reprise plus précisément par Joule en 1845).

\begin{enumerate}
\item Justifier que $W = 0$, puis que $Q = 0$. En déduire $\Delta U$, puis la température finale $T_f$ pour ce gaz parfait.
\item Pourquoi est-il impossible de représenter cette transformation par un chemin continu dans le diagramme de Clapeyron $(P,V)$ ?
\item On réalise, à partir du même état initial, une détente \emph{isotherme quasi-statique} (piston déplacé lentement) vers l'état $(T_0, 2V_0)$. Calculer le travail $W'$ et le transfert thermique $Q'$ pour cette seconde transformation.
\item Comparer les résultats des questions 1 et 3, alors que l'état final (même température $T_0$, même volume $2V_0$) est identique dans les deux cas. Qu'en conclure sur la nature de $W$ et de $Q$ ?
\item Peut-on conclure de $Q = 0$ que la détente de la question 1 ne comporte aucune irréversibilité ? Sur quel argument simple, indépendant de tout calcul, peut-on déjà affirmer le contraire ?
\end{enumerate}

\begin{indication}
Le gaz se détend dans le vide : aucune paroi mobile à pousser, $W = 0$ sans qu'il soit besoin d'aucune formule. Une transformation quasi-statique exige l'équilibre à chaque instant intermédiaire. Pour la question 5, penser au caractère spontané (ou non) du processus inverse.
\end{indication}

\begin{solution}
\textbf{Question 1.} $W = 0$ (pas de paroi à pousser), $Q = 0$ (parois adiabatiques). Donc $\Delta U = 0$, et pour un gaz parfait dont l'énergie interne ne dépend que de $T$, cela impose $T_f = T_0$.

\textbf{Question 2.} Pendant la détente, le système n'est à aucun instant à l'équilibre : la pression n'y est pas définie uniformément (le gaz se rue dans le vide de façon chaotique). Seuls les états initial et final, une fois l'équilibre revenu, sont représentables dans le diagramme de Clapeyron ; aucun chemin continu ne peut les relier.

\textbf{Question 3.} Détente isotherme quasi-statique de $V_0$ à $2V_0$, à $T_0$ constante :
\[
W' = -nRT_0\ln 2 < 0, \qquad Q' = -W' = nRT_0\ln 2 > 0
\]
(le gaz reçoit un travail négatif -- il en fournit -- et doit recevoir de la chaleur pour maintenir $T_0$ constante malgré la détente).

\textbf{Question 4.} Les états initial et final sont identiques dans les deux scénarios ($\Delta U = 0$ dans les deux cas, puisque $U$ ne dépend que de l'état via $T$), mais $W = Q = 0$ d'un côté, tandis que $W' \neq 0$ et $Q' \neq 0$ de l'autre. Contrairement à $\Delta U$, qui ne dépend que des états extrêmes, $W$ et $Q$ dépendent donc du \emph{chemin} suivi entre ces deux états : $\delta W$ et $\delta Q$ ne sont pas des différentielles exactes.

\textbf{Question 5.} Non : $Q = 0$ signifie seulement qu'aucune chaleur n'est \emph{échangée avec l'extérieur}, pas que la transformation soit exempte d'irréversibilité interne. Le signe de cette irréversibilité est déjà visible sans aucun calcul : on n'observe jamais le gaz occupant $2V_0$ se reconcentrer spontanément dans le compartiment de gauche, alors que rien dans le premier principe seul ne l'interdirait (l'énergie totale resterait inchangée). La leçon 5 permettra de quantifier précisément cette irréversibilité à l'aide de l'entropie -- \emph{qui n'est pas nulle bien que $Q$ le soit}, un résultat en apparence paradoxal traité en détail dans l'exercice « Détente de Joule et l'astuce fondamentale ».
\end{solution}
\end{exo}

\begin{exo}[Deux gaz séparés par une paroi interne]{equilibre-deux-gaz-paroi}
\keywords{premier principe, système isolé, paroi diatherme, paroi mobile, équilibre thermique, équilibre mécanique}

Un récipient rigide, isolé de l'extérieur, contient deux gaz parfaits monoatomiques séparés par une paroi interne initialement fixe, étanche et adiabatique. Le compartiment 1 contient $n_1$ moles à la température $T_1$ ; le compartiment 2 contient $n_2 = 4n_1$ moles à $T_2 = T_1/2$. Les deux compartiments ont le même volume $V$.

\begin{enumerate}
\item Exprimer les pressions initiales $P_1$ et $P_2$ et montrer que $P_2 = 2P_1$.
\item À $t=0$, la paroi devient diatherme (mais reste fixe et étanche). En appliquant le premier principe au système global, déterminer la température finale $T_f$ et les pressions finales de chaque compartiment, en fonction de $T_1$ et $P_1$.
\item Quelle énergie a été échangée entre les deux compartiments (en fonction de $n_1$, $R$, $T_1$) ? Sous quelle forme ?
\item On suppose maintenant qu'à $t=0$ la paroi devient diatherme \emph{et} mobile. Déterminer la température, la pression commune et les volumes finaux. Comparer les variations d'énergie interne de chaque gaz à celles de la question 2 et commenter la forme des échanges d'énergie.
\item Application numérique pour $n_1 = 1$ mol, $T_1 = 600$ K.
\end{enumerate}

\begin{indication}
Le système global est isolé (parois externes rigides et adiabatiques) : $\Delta U_1 + \Delta U_2 = 0$ dans tous les cas. Pour la question 4, l'équilibre impose l'égalité des températures \emph{et} des pressions.
\end{indication}

\begin{solution}
\textbf{Question 1.} $P_1 = \dfrac{n_1RT_1}{V}$ et $P_2 = \dfrac{4n_1R\,(T_1/2)}{V} = 2P_1$.

\textbf{Question 2.} Le système global n'échange ni travail (parois externes rigides) ni chaleur (parois adiabatiques) : $\Delta U_1 + \Delta U_2 = 0$, soit
\[
\tfrac{3}{2}n_1R(T_f - T_1) + \tfrac{3}{2}\,4n_1R\!\left(T_f - \tfrac{T_1}{2}\right) = 0 \implies 5T_f = 3T_1 \implies T_f = \tfrac{3}{5}T_1.
\]
Les volumes n'ont pas changé, donc $P_1^f = P_1\,\dfrac{T_f}{T_1} = \tfrac{3}{5}P_1$ et $P_2^f = 2P_1\,\dfrac{T_f}{T_2} = \tfrac{12}{5}P_1$.

\textbf{Question 3.} Le gaz 2 (le plus froid) reçoit
\[
Q = \Delta U_2 = 6n_1R\!\left(\tfrac{3}{5}T_1 - \tfrac{1}{2}T_1\right) = \tfrac{3}{5}n_1RT_1,
\]
et le gaz 1 cède exactement l'opposé. La paroi étant fixe, aucun travail n'est échangé : ce transfert se fait intégralement sous forme de chaleur.

\textbf{Question 4.} Le bilan global est inchangé : $\Delta U_1 + \Delta U_2 = 0$, et comme $U$ ne dépend que de $T$ pour un gaz parfait, on retrouve \emph{la même température finale} $T_f = \tfrac{3}{5}T_1$. L'équilibre mécanique impose de plus une pression commune :
\[
P = \frac{(n_1+n_2)RT_f}{2V} = \frac{5n_1R\cdot\tfrac{3}{5}T_1}{2V} = \tfrac{3}{2}P_1,
\]
avec $V_1^f = \dfrac{n_1RT_f}{P} = \tfrac{2V}{5}$ et $V_2^f = \tfrac{8V}{5}$ : la paroi s'est déplacée vers le gaz 1. Les variations d'énergie interne sont \emph{identiques} à celles de la question 2 ($\Delta U_i = \tfrac{3}{2}n_iR\,\Delta T_i$, mêmes températures extrêmes), mais cette fois l'énergie s'est échangée à la fois sous forme de chaleur et de travail : seul le bilan $Q + W = \Delta U$ est contraint, pas la répartition.

\textbf{Question 5.} $T_f = 360$ K et $|Q| = \tfrac{3}{5}\times 8{,}314 \times 600 \approx 3{,}0$ kJ.
\end{solution}
\end{exo}

\begin{exo}[Adiabatique réversible d'un gaz parfait : lois de Laplace]{lois-de-laplace}
\keywords{loi de Laplace, adiabatique réversible, indice adiabatique, gamma, détente adiabatique, premier principe différentiel}

On considère $n$ moles de gaz parfait d'indice adiabatique $\gamma = C_P/C_V$ supposé constant. On admet la relation de Mayer $C_P - C_V = nR$ (voir l'exercice sur la découverte de $C_P$).

\begin{enumerate}
\item Montrer que
\[
C_V = \frac{nR}{\gamma - 1}, \qquad C_P = \frac{\gamma\, nR}{\gamma - 1}.
\]
\item Le gaz subit une transformation adiabatique réversible. En partant du premier principe sous forme différentielle, montrer que
\[
\frac{\mathrm{d}T}{T} + (\gamma - 1)\frac{\mathrm{d}V}{V} = 0,
\]
puis en déduire que $TV^{\gamma-1} = \text{cste}$ le long de la transformation.
\item En déduire les deux autres formes de la loi de Laplace : $PV^\gamma = \text{cste}$ et $T^\gamma P^{1-\gamma} = \text{cste}$.
\item Détente adiabatique réversible qui double le volume, pour un gaz monoatomique ($\gamma = 5/3$), $n = 1$ mol, $T_I = 300$ K. Calculer $T_F$ puis $W$, et commenter les signes.
\item Dans le diagramme de Clapeyron, comparer en un même point $(V,P)$ la pente d'une adiabatique réversible et celle d'une isotherme.
\end{enumerate}

\begin{indication}
Pour la question 2 : $\mathrm{d}U = C_V\,\mathrm{d}T$ et $\delta W = -P\,\mathrm{d}V$ avec $P = nRT/V$, puis séparer les variables. Pour la question 5, différencier $PV^\gamma = \text{cste}$ et $PV = \text{cste}$.
\end{indication}

\begin{solution}
\textbf{Question 1.} De $C_P = \gamma C_V$ et $C_P - C_V = nR$ : $C_V(\gamma - 1) = nR$, d'où les deux formules.

\textbf{Question 2.} Adiabatique : $\delta Q = 0$, donc $\mathrm{d}U = \delta W$, soit $C_V\,\mathrm{d}T = -P\,\mathrm{d}V = -\dfrac{nRT}{V}\mathrm{d}V$. Avec $C_V = nR/(\gamma-1)$ :
\[
\frac{\mathrm{d}T}{T} = -(\gamma-1)\frac{\mathrm{d}V}{V}.
\]
En intégrant : $\ln T + (\gamma-1)\ln V = \text{cste}$, c'est-à-dire $TV^{\gamma-1} = \text{cste}$.

\textbf{Question 3.} En substituant $T = PV/(nR)$ : $PV^\gamma = \text{cste}$. En substituant $V = nRT/P$ : $T^\gamma P^{1-\gamma} = \text{cste}$.

\textbf{Question 4.} $T_F = T_I\,2^{-(\gamma-1)} = 300 \times 2^{-2/3} \approx 189$ K. Comme $Q = 0$ :
\[
W = \Delta U = \tfrac{3}{2}R(T_F - T_I) \approx \tfrac{3}{2}\times 8{,}314 \times (-111) \approx -1{,}4\ \text{kJ}.
\]
$W < 0$ : le gaz fournit du travail à l'extérieur. Ne pouvant le prélever sur aucune source de chaleur ($Q=0$), il puise dans son énergie interne : d'où le refroidissement $\Delta T < 0$.

\textbf{Question 5.} En différenciant $PV = \text{cste}$ : $\left(\frac{\mathrm{d}P}{\mathrm{d}V}\right)_{\text{isoT}} = -\dfrac{P}{V}$. En différenciant $PV^\gamma = \text{cste}$ : $\left(\frac{\mathrm{d}P}{\mathrm{d}V}\right)_{\text{adiab}} = -\gamma\dfrac{P}{V}$. L'adiabatique est $\gamma$ fois plus pentue que l'isotherme : comprimer sans évacuer la chaleur fait monter la pression plus vite, car la température augmente aussi.
\end{solution}
\end{exo}

\begin{exo}[Travail d'une transformation polytropique $PV^k=\text{cste}$]{transformation-polytropique}
\keywords{transformation polytropique, travail des forces de pression, adiabatique reversible, loi de Laplace}

On considere une transformation quasi-statique d'un gaz parfait telle que $PV^k=\text{cste}$, $k\in\mathbb{R}$, $k\neq1$.

\begin{enumerate}
\item Calculer $W_{A\to B}$ en fonction de $P_A,V_A,P_B,V_B$ et $k$.
\item Retrouver le resultat pour une isobare ($k=0$).
\item Conclure directement pour une isochore sans utiliser la formule generale.
\item Expliquer pourquoi $k=1$ est exclu, puis retrouver $W=-nRT\ln(V_B/V_A)$.
\item Pour $k=\gamma=5/3$ (adiabatique reversible, monoatomique), montrer que $W=\Delta U$.
\end{enumerate}

\begin{indication}
Poser $P(V)=CV^{-k}$ avec $C=P_AV_A^k$, integrer, puis utiliser $CV^{1-k}=PV$.
\end{indication}

\begin{solution}
\textbf{Question 1.}
\[
W = \frac{P_AV_A - P_BV_B}{1-k}
\]

\textbf{Question 2.} $k=0$, $P_A=P_B=P$ : $W=P(V_A-V_B)=-P\Delta V$. \checkmark

\textbf{Question 3.} Isochore : $\mathrm{d}V=0$ a chaque instant, donc $W=0$. La formule polytropique ne s'applique pas ($k\to\infty$).

\textbf{Question 4.} Pour $k=1$ le facteur $1/(1-k)$ diverge (primitive logarithmique). Avec $PV=nRT=\text{cste}$ :
\[
W = -nRT\ln\!\tfrac{V_B}{V_A}.
\]

\textbf{Question 5.} $k=5/3$ :
\[
W = \tfrac{3}{2}(P_BV_B-P_AV_A) = \tfrac{3}{2}nR(T_B-T_A) = \Delta U.
\]
Coherent avec le premier principe ($Q=0$). \checkmark
\end{solution}
\end{exo}

\begin{exo}[Compression adiabatique : progressive ou brutale]{compression-progressive-brutale}
\keywords{compression adiabatique, transformation réversible, transformation irréversible, loi de Laplace, pression extérieure constante}

Un cylindre vertical de section $S = 200\ \text{cm}^2$, aux parois calorifugées, est fermé par un piston calorifugé de masse négligeable, mobile sans frottement. Il contient $m = 2{,}0$ g d'hélium (gaz parfait monoatomique, $M = 4{,}0$ g/mol, $\gamma = 5/3$) dans les conditions initiales $T_0 = 300$ K, $P_0 = 1{,}0$ bar. On prendra $g = 10\ \text{m·s}^{-2}$.

\begin{enumerate}
\item Calculer le nombre de moles $n$ et le volume initial $V_0$.
\item On dépose très progressivement, grain par grain, une masse totale $M_s = 100$ kg de sable sur le piston. Justifier que la transformation peut être considérée comme réversible, et calculer le taux de compression $x = P_1/P_0$. La valeur de $x$ dépend-elle de la façon de déposer le sable ?
\item Calculer $T_1$ et $V_1$ dans l'état final.
\item On repart de l'état initial et on dépose cette fois les 100 kg d'un seul coup : le gaz subit brutalement la pression extérieure constante $P_1$. En appliquant le premier principe, montrer que
\[
T_2 = T_0\,\frac{1 + (\gamma-1)\,x}{\gamma},
\]
et calculer $T_2$ et $V_2$.
\item Comparer $T_1$ et $T_2$. Laquelle des deux compressions échauffe le plus le gaz ? Interpréter.
\item Depuis l'état de la question 4, on retire brutalement tout le sable. Déterminer l'état final $(T_3, V_3)$. Revient-on à l'état initial ? Conclure.
\end{enumerate}

\begin{indication}
Question 2 : équilibre des forces sur le piston. Question 4 : $W = -P_1(V_2 - V_0)$ car la pression extérieure est constante, et à l'équilibre final $P_2 = P_1$ ; éliminer les volumes avec l'équation d'état. Question 6 : refaire le même raisonnement avec $P_{\text{ext}} = P_0$ depuis $(T_2, P_1)$.
\end{indication}

\begin{solution}
\textbf{Question 1.} $n = m/M = 0{,}50$ mol et $V_0 = nRT_0/P_0 \approx 12{,}5$ L.

\textbf{Question 2.} Chaque grain ne déséquilibre le piston que de façon infinitésimale : le gaz passe par une succession continue d'états d'équilibre (transformation quasi-statique sans frottement, donc réversible). L'équilibre final du piston donne
\[
P_1 = P_0 + \frac{M_s\,g}{S} = 10^5 + \frac{1000}{2\times10^{-2}} = 1{,}5\times10^5\ \text{Pa}, \qquad x = 1{,}5.
\]
$x$ ne dépend que de l'équilibre des forces, pas du chemin : il serait le même en déposant la masse d'un coup.

\textbf{Question 3.} Loi de Laplace $T^\gamma P^{1-\gamma} = \text{cste}$ :
\[
T_1 = T_0\,x^{(\gamma-1)/\gamma} = 300\times 1{,}5^{2/5} \approx 353\ \text{K}, \qquad V_1 = V_0\,x^{-1/\gamma} \approx 9{,}8\ \text{L}.
\]

\textbf{Question 4.} La transformation est brutale : seule la pression extérieure $P_1$ est définie, et $W = -P_1(V_2 - V_0)$, $Q = 0$. Premier principe :
\[
\tfrac{3}{2}nR(T_2 - T_0) = -P_1\!\left(\frac{nRT_2}{P_1} - \frac{nRT_0}{P_0}\right) = -nRT_2 + x\,nRT_0.
\]
D'où $\tfrac{5}{2}T_2 = \left(\tfrac{3}{2} + x\right)T_0$, ce qui est bien $T_2 = T_0\,\dfrac{1 + (\gamma-1)x}{\gamma}$ pour $\gamma = 5/3$. Numériquement $T_2 = \tfrac{6}{5}T_0 = 360$ K et $V_2 = nRT_2/P_1 \approx 10{,}0$ L.

\textbf{Question 5.} $T_2 = 360\ \text{K} > T_1 \approx 353\ \text{K}$ : la compression brutale échauffe davantage, à taux de compression identique. Pendant la transformation brutale, le gaz est hors équilibre et la totalité du travail $-P_1\Delta V$ (calculé avec la pression finale, la plus élevée) lui est fournie : il reçoit plus de travail que dans le cas réversible. Cette différence est la signature de l'irréversibilité, que le second principe quantifiera à la leçon suivante.

\textbf{Question 6.} Même raisonnement depuis $(T_2, P_1)$ avec $P_{\text{ext}} = P_0$, soit un rapport $P_0/P_1 = 2/3$ :
\[
T_3 = T_2\,\frac{1 + (\gamma-1)\cdot\tfrac{2}{3}}{\gamma} = 360\times\frac{13}{15} = 312\ \text{K}, \qquad V_3 = \frac{nRT_3}{P_0} \approx 13{,}0\ \text{L}.
\]
On ne revient pas à l'état initial : après un aller-retour brutal, le gaz est plus chaud ($312$ K au lieu de $300$ K) et plus dilaté. Le cycle de manipulations est identique, mais l'histoire ne s'efface pas : les transformations brutales sont irréversibles.
\end{solution}
\end{exo}

\begin{exo}[Cycle biisobare-biisochore : bilan energetique complet]{cycle-biisobare-biisochore}
\keywords{cycle thermodynamique, diagramme de Clapeyron, moteur thermique, gaz parfait}

On considere $n$ moles de gaz parfait monoatomique ($U=\frac{3}{2}nRT$) decrivant, dans le plan de Clapeyron $(V,P)$, le cycle a deux isochores et deux isobares, parcouru dans le sens horaire :
\[
1(V_1,P_1)\ \xrightarrow{\text{isochore}}\ 3(V_1,P_2)\ \xrightarrow{\text{isobare}}\ 2(V_2,P_2)\ \xrightarrow{\text{isochore}}\ 4(V_2,P_1)\ \xrightarrow{\text{isobare}}\ 1(V_1,P_1)
\]
avec $P_2>P_1$ et $V_2>V_1$.

\begin{enumerate}
\item Exprimer $\Delta U$, $W$ et $Q$ sur chacune des quatre branches, en fonction de $P_1,P_2,V_1,V_2$ uniquement.
\item Verifier que $\Delta U_{\text{total}}=0$ et que $Q_{\text{total}}=-W_{\text{total}}$.
\item Application numerique : $P_1=1{,}0$ bar, $P_2=3{,}0$ bar, $V_1=5{,}0$ L, $V_2=15{,}0$ L.
\item Sur quelles branches le systeme recoit-il de la chaleur ? Commenter.
\end{enumerate}

\begin{indication}
Utiliser $PV=nRT$ pour eliminer $T$ : sur une isochore $nR\,\Delta T = V\,\Delta P$, sur une isobare $nR\,\Delta T = P\,\Delta V$. Le travail elementaire vaut $\delta W = -P_{\text{ext}}\mathrm{d}V = -P\,\mathrm{d}V$ (quasi-statique). Convention banquier.
\end{indication}

\begin{solution}
\textbf{Branche $1\to3$ (isochore, $V=V_1$).}
\[
W=0,\qquad \Delta U = \tfrac{3}{2}(P_2-P_1)V_1,\qquad Q=\Delta U
\]

\textbf{Branche $3\to2$ (isobare, $P=P_2$).}
\[
W=-P_2(V_2-V_1),\qquad \Delta U=\tfrac{3}{2}P_2(V_2-V_1),\qquad Q=\tfrac{5}{2}P_2(V_2-V_1)
\]

\textbf{Branche $2\to4$ (isochore, $V=V_2$).}
\[
W=0,\qquad \Delta U=-\tfrac{3}{2}(P_2-P_1)V_2,\qquad Q=\Delta U
\]

\textbf{Branche $4\to1$ (isobare, $P=P_1$).}
\[
W=P_1(V_2-V_1),\qquad \Delta U=-\tfrac{3}{2}P_1(V_2-V_1),\qquad Q=-\tfrac{5}{2}P_1(V_2-V_1)
\]

En sommant : $\Delta U_{\text{total}}=0$ et $W_{\text{total}}=-(P_2-P_1)(V_2-V_1)<0$, $Q_{\text{total}}=(P_2-P_1)(V_2-V_1)>0$.

\textbf{Application numerique} ($\Delta P=2{,}0\times10^5$ Pa, $\Delta V=1{,}0\times10^{-2}$ m$^3$) :
\begin{align*}
1\to3:\ & W=0,\ Q=1500\text{ J}\\
3\to2:\ & W=-3000\text{ J},\ Q=7500\text{ J}\\
2\to4:\ & W=0,\ Q=-4500\text{ J}\\
4\to1:\ & W=1000\text{ J},\ Q=-2500\text{ J}
\end{align*}
Total : $\Delta U=0$, $W=-2000$ J, $Q=+2000$ J.

Le systeme recoit de la chaleur sur les branches $1\to3$ et $3\to2$ (montee en pression puis detente a haute pression). Sur le reste du cycle, il restitue de la chaleur.
\end{solution}
\end{exo}

\begin{exo}[Cycle triangulaire]{cycle-triangulaire-diagonale}
\keywords{cycle moteur, transformation lineaire P(V), diagramme de Clapeyron, aire du cycle}

On considere $n$ moles de gaz parfait monoatomique decrivant le cycle $A\to C\to B\to A$, avec
\[
A(V_0,P_0),\quad C(V_0,2P_0),\quad B(2V_0,P_0),
\]
$A\to C$ isochore, $B\to A$ isobare, $C\to B$ segment de droite dans le plan $(V,P)$.

\begin{enumerate}
\item Verifier que le sens $A\to C\to B\to A$ est bien horaire.
\item Determiner $P(V)$ de la droite $C\to B$.
\item Calculer $\Delta U$, $W$ et $Q$ sur chacune des trois branches.
\item Verifier $\Delta U_{\text{total}}=0$ et que $W_{\text{total}}$ egale l'aire du triangle.
\item Montrer que $\Delta U_{C\to B}=0$ bien que la transformation ne soit pas isotherme.
\end{enumerate}

\begin{indication}
Pour $C\to B$, chercher $P(V)=aV+b$ en imposant le passage par $C$ et $B$. Pour la question 5, $\Delta U$ ne depend que des etats initial et final.
\end{indication}

\begin{solution}
\textbf{Questions 1--2.} Cycle horaire (moteur). Droite $C\to B$ : $P(V) = P_0(3-V/V_0)$.

\textbf{Branche $A\to C$.} $W=0$, $\Delta U=\tfrac32P_0V_0$, $Q=\tfrac32P_0V_0$.

\textbf{Branche $C\to B$.}
\[
W = -\int_{V_0}^{2V_0}P_0\!\left(3-\tfrac{V}{V_0}\right)\mathrm{d}V = -\tfrac{3}{2}P_0V_0.
\]
$P_CV_C=2P_0V_0=P_BV_B \Rightarrow T_B=T_C \Rightarrow \Delta U=0$, donc $Q=\tfrac32P_0V_0$.

\textbf{Branche $B\to A$.} $W=P_0V_0$, $\Delta U=-\tfrac32P_0V_0$, $Q=-\tfrac52P_0V_0$.

\textbf{Questions 4--5.} $\Delta U_{\text{total}}=0$, $W_{\text{total}}=-\tfrac12P_0V_0=$ aire du triangle. $\Delta U_{C\to B}=0$ car $U$ est une fonction d'etat : seuls les etats extremes comptent, et ici $T_C=T_B$ par coincidence geometrique.
\end{solution}
\end{exo}

\begin{exo}[Cycle de Lenoir]{cycle-lenoir}
\keywords{cycle de Lenoir, isobare, isotherme, isochore, diagramme de Clapeyron, cycle récepteur, gaz diatomique}

Une mole de gaz parfait diatomique ($C_V = \frac{5}{2}R$, $C_P = \frac{7}{2}R$) est initialement dans l'état $A$ de coordonnées $P_0 = 2{,}0\times10^5$ Pa et $V_0 = 14$ L dans le diagramme de Clapeyron. On lui fait subir successivement :
\begin{itemize}
\item une détente isobare $A\to B$ qui triple son volume ;
\item une compression isotherme quasi-statique $B\to C$ qui le ramène au volume initial ;
\item un refroidissement isochore $C\to A$ qui le ramène à l'état initial.
\end{itemize}

\begin{enumerate}
\item Représenter le cycle $ABCA$ dans le diagramme de Clapeyron. Calculer $T_A$, la température $T_B$ de l'isotherme, et la pression maximale $P_C$ atteinte.
\item Calculer $W$ et $Q$ sur chacune des trois branches, puis sur le cycle complet. On remarquera que $RT_A = P_0V_0$.
\item Le cycle ainsi parcouru est-il moteur ou récepteur ? Retrouver ce résultat à partir du sens de parcours dans le diagramme. Comment faudrait-il le parcourir pour en faire un moteur ?
\end{enumerate}

\begin{indication}
$T_A = P_0V_0/R$ et sur l'isobare $T \propto V$. Sur l'isotherme, $\Delta U = 0$ et $W = -nRT\ln(V_f/V_i)$. Le sens antihoraire correspond à un cycle récepteur.
\end{indication}

\begin{solution}
\textbf{Question 1.} $T_A = \dfrac{P_0V_0}{R} = \dfrac{2\times10^5 \times 14\times10^{-3}}{8{,}314} \approx 337$ K. Sur l'isobare, $T_B = 3T_A \approx 1010$ K. Sur l'isotherme $B\to C$ : $P_CV_0 = RT_B = 3P_0V_0$, donc $P_C = 3P_0 = 6{,}0\times10^5$ Pa. Le cycle est bordé par l'isobare $P_0$ en bas, l'isotherme $T_B$ en haut, et l'isochore $V_0$ à gauche.

\textbf{Question 2.} On pose $P_0V_0 = RT_A = 2800$ J.

$A\to B$ (isobare, $\Delta V = 2V_0$) :
\[
W_{AB} = -P_0\cdot 2V_0 = -5600\ \text{J},\quad \Delta U_{AB} = \tfrac{5}{2}R\cdot 2T_A = 5RT_A = 14\,000\ \text{J},\quad Q_{AB} = \tfrac{7}{2}R\cdot2T_A = 19\,600\ \text{J}.
\]

$B\to C$ (isotherme à $3T_A$) : $\Delta U = 0$ et
\[
W_{BC} = -3RT_A\ln\frac{V_0}{3V_0} = 3RT_A\ln 3 \approx +9230\ \text{J}, \qquad Q_{BC} = -W_{BC} \approx -9230\ \text{J}.
\]

$C\to A$ (isochore) : $W_{CA} = 0$ et $Q_{CA} = \Delta U_{CA} = \tfrac{5}{2}R(T_A - 3T_A) = -5RT_A = -14\,000$ J.

Cycle complet : $\Delta U = 0$, $W_{\text{total}} = -5600 + 9230 \approx +3630$ J et $Q_{\text{total}} \approx -3630$ J.

\textbf{Question 3.} $W_{\text{total}} > 0$ : le gaz reçoit du travail et rejette de la chaleur, le cycle est récepteur. C'est cohérent avec le sens de parcours : on parcourt l'isobare (branche basse) de gauche à droite et l'isotherme (branche haute) de droite à gauche, soit un parcours antihoraire. Parcouru dans l'autre sens ($A\to C\to B\to A$), le même cycle fournirait $3630$ J de travail par cycle : c'est le principe du moteur de Lenoir, l'un des premiers moteurs à combustion interne (1860).
\end{solution}
\end{exo}

\begin{exo}[Cycle frigorifique : le meme cycle parcouru a l'envers]{cycle-frigorifique-cop}
\keywords{cycle recepteur, refrigerateur, coefficient de performance, machine thermique}

On reprend le cycle biisobare-biisochore de l'exercice \textit{cycle-biisobare-biisochore} (memes valeurs numeriques), mais parcouru dans le sens \textbf{antihoraire} : $1\to4\to2\to3\to1$.

\begin{enumerate}
\item Justifier que chaque $W$ et $Q$ change de signe. Nature de ce cycle ?
\item Identifier les branches ou le systeme absorbe / cede de la chaleur.
\item Calculer $Q_c$, $|Q_h|$ et $W$.
\item En deduire le $\text{COP}=Q_c/W$. Commenter.
\end{enumerate}

\begin{indication}
Parcourir un cycle a l'envers change le signe de $W$ et $Q$ sur chaque branche. Reutiliser les resultats numeriques de l'exercice precedent.
\end{indication}

\begin{solution}
\textbf{Question 1.} Le cycle antihoraire est recepteur : $W_{\text{total}}=+2000$ J recu, $Q_{\text{total}}=-2000$ J cede. C'est le principe d'un refrigerateur.

\textbf{Question 2.}
\begin{align*}
1\to4 &: Q=+2500\text{ J (absorbee, source froide)}\\
4\to2 &: Q=+4500\text{ J (absorbee, source froide)}\\
2\to3 &: Q=-7500\text{ J (cedee, source chaude)}\\
3\to1 &: Q=-1500\text{ J (cedee, source chaude)}
\end{align*}

\textbf{Question 3.} $Q_c=7000$ J, $|Q_h|=9000$ J, $W=2000$ J. Verification : $7000-9000+2000=0$. \checkmark

\textbf{Question 4.} $\text{COP}=7000/2000=3{,}5$. Pour chaque joule fourni au compresseur, $3{,}5$ J sont extraits de la source froide --- un COP superieur a 1 est normal pour un refrigerateur, contrairement au rendement d'un moteur thermique.
\end{solution}
\end{exo}

\begin{exo}[Résistance chauffante dans un cylindre cloisonné]{resistance-chauffante-cloison}
\keywords{premier principe, piston adiabatique, loi de Laplace, compression adiabatique réversible, bilan d'énergie, hélium}

Un cylindre horizontal rigide, aux parois calorifugées, est divisé en deux compartiments $C_1$ et $C_2$ par un piston mobile sans frottement, lui aussi calorifugé. Chaque compartiment contient de l'hélium (gaz parfait monoatomique, $\gamma = 5/3$) dans le même état initial : $P_0 = 1{,}0\times10^5$ Pa, $V_0 = 1{,}0$ L, $T_0 = 300$ K. Une résistance électrique placée dans $C_1$ lui apporte de la chaleur très lentement, jusqu'à ce que la pression commune atteigne $P_1 = 1{,}5\times10^5$ Pa. On supposera toutes les transformations quasi-statiques.

\begin{enumerate}
\item Quelle transformation subit le gaz de $C_2$ ? En déduire son volume $V_2$ et sa température $T_2$ finaux.
\item En déduire le volume $V_1$ et la température $T_1$ finaux du gaz de $C_1$.
\item Calculer les variations d'énergie interne $\Delta U_1$ et $\Delta U_2$. On remarquera que pour un gaz parfait monoatomique, $U = \frac{3}{2}PV$.
\item En appliquant le premier principe au système global $\{C_1 + C_2\}$, déterminer la chaleur $Q$ fournie par la résistance. Vérifier la cohérence en appliquant le premier principe à $C_2$ seul, puis à $C_1$ seul.
\end{enumerate}

\begin{indication}
Le gaz de $C_2$ n'échange aucune chaleur (piston et parois calorifugés) et est comprimé de façon quasi-statique : c'est une adiabatique réversible, la loi de Laplace $PV^\gamma = \text{cste}$ s'applique. Le volume total $V_1 + V_2 = 2V_0$ est fixé. Pour le système global, les parois externes sont rigides : seul $Q$ entre.
\end{indication}

\begin{solution}
\textbf{Question 1.} Le gaz de $C_2$ est comprimé lentement, sans frottement et sans échange de chaleur : il subit une compression adiabatique réversible. Loi de Laplace :
\[
V_2 = V_0\!\left(\frac{P_0}{P_1}\right)^{1/\gamma} = 1{,}0\times\left(\frac{1}{1{,}5}\right)^{3/5} \approx 0{,}78\ \text{L},
\qquad
T_2 = T_0\,\frac{P_1V_2}{P_0V_0} \approx 353\ \text{K}.
\]

\textbf{Question 2.} Le cylindre est rigide : $V_1 = 2V_0 - V_2 \approx 1{,}22$ L, et l'équation d'état donne
\[
T_1 = T_0\,\frac{P_1V_1}{P_0V_0} = 300 \times 1{,}5 \times 1{,}22 \approx 547\ \text{K}.
\]

\textbf{Question 3.} Avec $U = \tfrac{3}{2}PV$ :
\[
\Delta U_1 = \tfrac{3}{2}(P_1V_1 - P_0V_0) \approx \tfrac{3}{2}(182{,}4 - 100) \approx 124\ \text{J},
\qquad
\Delta U_2 = \tfrac{3}{2}(P_1V_2 - P_0V_0) \approx \tfrac{3}{2}(117{,}6 - 100) \approx 26\ \text{J}.
\]

\textbf{Question 4.} Le système global n'échange aucun travail (parois externes rigides) et reçoit uniquement $Q$ :
\[
Q = \Delta U_1 + \Delta U_2 = \tfrac{3}{2}(P_1 - P_0)(2V_0) = \tfrac{3}{2}\times 0{,}5\times10^5 \times 2\times10^{-3} = 150\ \text{J}.
\]
Cohérence : $C_2$ est adiabatique, donc $W_2 = \Delta U_2 \approx 26$ J est le travail que lui fournit le piston. Le gaz de $C_1$ reçoit $Q$ de la résistance et cède ce travail au piston : $\Delta U_1 = Q - W_2 \approx 150 - 26 = 124$ J. \checkmark

On note l'efficacité du bilan global : nul besoin de connaître le détail des échanges internes pour trouver $Q$.
\end{solution}
\end{exo}
